\documentclass[12pt]{article}
\usepackage[utf8]{inputenc}
\usepackage[T1]{fontenc}
\usepackage[margin=1in]{geometry}
\usepackage{setspace}
\usepackage{xcolor}
\usepackage[colorlinks=true, urlcolor=blue!60!black, linkcolor=blue!60!black]{hyperref}
\setlength{\parindent}{2em}
\setlength{\parskip}{1.25em}
\renewcommand{\baselinestretch}{1.0}
\begin{document}
\begin{center}
{\LARGE \textit{Miller v. California}, 413 U.S. 15 (1973)\par}
\vspace{0.8em}
{\large Case Brief --- Four-Page Edition\par}
\end{center}
\vspace{1em}

\section{Memory Jogger}

The Supreme Court reaffirmed that obscene speech lies outside First Amendment protection and replaced the ``utterly without redeeming social value'' test with a three-prong obscenity framework anchored in contemporary local community standards and lacking serious literary, artistic, political, or scientific value (\href{https://supreme.justia.com/cases/federal/us/413/15/}{413 U.S. 15}).

\section{Detailed Case Facts}

Marvin Miller ran a California mail-order business in adult books and films. His prosecution arose from five unsolicited brochures mailed to a restaurant in Newport Beach; the manager and his mother, who had never requested them, complained to police. The brochures advertised four books, including \textit{Sex Orgies Illustrated} and \textit{An Illustrated History of Pornography}, and a film, \textit{Marital Intercourse}, and consisted mainly of photographs and drawings explicitly depicting group sexual activity with genitals prominently displayed.

Miller was charged with a misdemeanor under California Penal Code \S~311.2(a) for knowingly distributing obscene matter. The statute then defined ``obscene'' largely along the lines of the \textit{Memoirs v. Massachusetts} plurality, requiring that the matter be ``utterly without redeeming social importance.''

\section{Procedural History}

At trial in Orange County, both sides assumed statewide standards governed, and the judge instructed the jury on the ``contemporary community standards of the State of California'' without objection. The jury convicted. On appeal, Miller argued for the first time that a national standard was required; the Appellate Department of the Superior Court summarily affirmed. The Supreme Court heard argument in 1972, ordered reargument, and decided the case on June 21, 1973, with \textit{Paris Adult Theatre I v. Slaton}.

\section{Judicial Votes}

Vacated and remanded, \textbf{5--4}.
\begin{itemize}
\item \textbf{Majority:} Chief Justice Burger, joined by Justices White, Blackmun, Powell, and Rehnquist.
\item \textbf{Concurrences:} None.
\item \textbf{Dissents:} Justice Douglas, alone; Justice Brennan, joined by Justices Stewart and Marshall.
\end{itemize}

\section{Holding}

Reaffirming \textit{Roth v. United States}, obscene material is categorically unprotected. States may regulate only works depicting or describing sexual conduct specifically defined by state law, judged by three conjunctive guidelines: (a) whether the average person, applying contemporary community standards, would find that the work, taken as a whole, appeals to the prurient interest; (b) whether it depicts or describes, in a patently offensive way, sexual conduct specifically defined by state law; and (c) whether it, taken as a whole, lacks serious literary, artistic, political, or scientific value. The \textit{Memoirs} test is rejected, and juries may apply local or statewide rather than national standards to prongs (a) and (b). Burger's ``plain examples'' of hard-core conduct are illustrative dicta.

\section{Analysis of Opinions}

\textbf{Majority (Burger).} Burger sought majority consensus on ``the intractable obscenity problem'' after summary reversals under \textit{Redrup} had made the Court ``an unreviewable board of censorship for the 50 States.'' Miller's brochures were thrust on unwilling recipients, implicating the state interest in protecting unconsenting adults and juveniles. The \textit{Memoirs} plurality had forced prosecutors to prove a negative, a burden that never commanded a majority and that even its author, Justice Brennan, had abandoned. The new test isolates ``hard core'' pornography, gives fair notice through specific statutory definitions, and preserves independent appellate review. A single national standard would be ``an exercise in futility'' in a diverse nation, so California's statewide standard sufficed. Burger distinguished the ``unfettered interchange of ideas'' from ``commercial exploitation of obscene material.''

\textbf{Douglas (dissenting).} ``Obscenity'' appears nowhere in the Constitution, and no Founding-era exception excluded it from ``the press''; censorship requires a constitutional amendment, not judicial improvisation. Applying new standards retroactively carried ``all the evils of an \textit{ex post facto} law,'' so due process required a prior civil adjudication that a work was obscene, and an ``offensive'' standard cut ``the very vitals out of the First Amendment.''

\textbf{Brennan, with Stewart and Marshall (dissenting).} Relying on his \textit{Paris Adult Theatre} dissent, Brennan concluded that no formulation could separate obscene from protected speech precisely enough. Reserving laws aimed at juveniles or unconsenting adults, he found \S~311.2(a) substantially overbroad.

\section{Examples: Future Applications}

\begin{enumerate}
\item \textbf{Same-Side (Unprotected): Unsolicited Pamphlets.} A distributor mails graphic pamphlets of actual intercourse, with no commentary, into a conservative county whose statute enumerates those acts. All three prongs are met.
\item \textbf{Same-Side (Unprotected): Token Voltaire.} A hard-core magazine adds one quotation from Voltaire. Footnote 7 says a token quotation ``will not constitutionally redeem an otherwise obscene publication.''
\item \textbf{Opposite-Side (Protected): Medical Textbook.} An obstetrics manual with graphic anatomical photographs has serious scientific value under prong (c), judged objectively under \textit{Pope v. Illinois}.
\item \textbf{Opposite-Side (Protected): Acclaimed Film.} A vague ban on ``lewd, indecent, or immoral performances'' is applied to a film with non-explicit nudity. Nudity is not obscene, the conduct is not specifically defined, and the film has artistic value.
\item \textbf{Fence-Sitter: Explicit Political Film Online.} An explicit streamed film includes a feminist political monologue and is prosecuted in a rural district. Local standards may meet prongs (a) and (b), raising the Internet problem in \textit{Ashcroft v. ACLU}, while prong (c) asks whether the monologue is serious or a pretextual wrapper.
\end{enumerate}

\section{Critique}

\textbf{Originalist and textualist views.} Defenders cite Founding-era and nineteenth-century restrictions as proof that obscenity lay outside ``the freedom of speech'' and praise \textit{Miller}'s federalism. Textualists echoing Douglas reply that the text has no obscenity carve-out and that the test is an unmoored invention; Justice Scalia later doubted that artistic value can be litigated, since ``there is no use litigating about it.''

\textbf{Progressive, feminist, and doctrinal views.} Civil libertarians fault \textit{Miller}'s vagueness and chilling effect on publishers and artists, magnified online, where local standards risk imposing the most restrictive forum's norms on national speech. Catharine MacKinnon argued from the opposite direction that obscenity law targets moral offense rather than pornography's real harm, the subordination of women.

\section{Key Quotations}

\begin{quote}
``The basic guidelines for the trier of fact must be: (a) whether `the average person, applying contemporary community standards' would find that the work, taken as a whole, appeals to the prurient interest, \textit{Kois v. Wisconsin}, supra, at 230, quoting \textit{Roth v. United States}, supra, at 489; (b) whether the work depicts or describes, in a patently offensive way, sexual conduct specifically defined by the applicable state law; and (c) whether the work, taken as a whole, lacks serious literary, artistic, political, or scientific value.''
\end{quote}
\noindent --- Chief Justice Burger, Opinion of the Court, 413 U.S. at 24.

\begin{quote}
``It is neither realistic nor constitutionally sound to read the First Amendment as requiring that the people of Maine or Mississippi accept public depiction of conduct found tolerable in Las Vegas, or New York City.''
\end{quote}
\noindent --- Chief Justice Burger, Opinion of the Court, 413 U.S. at 32.

\begin{quote}
``The First Amendment protects works which, taken as a whole, have serious literary, artistic, political, or scientific value, regardless of whether the government or a majority of the people approve of the ideas these works represent.''
\end{quote}
\noindent --- Chief Justice Burger, Opinion of the Court, 413 U.S. at 34.

\end{document}
